Hiermit versichere ich, dass ich die vorgelegte Arbeit selbst verfasst und keine anderen
als die angegebenen Quellen und Hilfsmittel benutzt habe. Die aus fremden Quellen
direkt oder indirekt übernommenen Texte, Gedankengänge, Konzepte, Grafiken usw.
habe ich als solche gekennzeichnet und mit vollständigen Verweisen auf die jeweilige
Quelle versehen.

Mir ist bewusst, dass die Nutzung mittels generativer KI erstellter Texte oder Inhalte
keine Garantie für deren Qualität gewährleistet und ich die Verantwortung trage, falls
es durch die Verwendung solcher Hilfsmittel zu fehlerhaften Inhalten, zu Verstößen
gegen das Datenschutzrecht, Urheberrecht oder zu wissenschaftlichem Fehlverhalten
(z.B. Plagiate) kommt.

Ich versichere außerdem:

\begin{itemize}
	\item[$\mathrlap{\circ}\!\times$]
	dass ich mich generativer KI-Tools lediglich als Hilfsmittel bedient habe und
	dass ich angegeben habe, welche generativen KI-Tools ich zu welchem Zweck
	und in welchem Umfang eingesetzt habe. Eine entsprechende Tabelle habe ich
	dem Anhang meiner Arbeit beigefügt.
	\item[$\mathrlap{\circ}\!\times$]
	dass ich keine generativen KI-Tools zur Erstellung von Texten oder sonstiger
	Inhalte verwendet habe.
\end{itemize}

Ich versichere weiterhin, dass die Arbeit weder vollständig noch in wesentlichen Teilen
Gegenstand eines anderen Prüfungsverfahrens gewesen ist, dass ich die Arbeit weder
vollständig noch in wesentlichen Teilen bereits veröffentlicht habe, dass die Inhalte der
für die Plagiatsprüfung eingereichten Datei mit den Inhalten des eingereichten
gedruckten Berichts übereinstimmen, und ich mir bewusst bin, dass eine
Plagiatsprüfung durchgeführt wird.

Mir ist bekannt, dass ein Verstoß gegen die genannten Punkte prüfungsrechtliche
Konsequenzen haben und insbesondere dazu führen kann, dass die Prüfungsleistung
mit „nicht ausreichend“ bzw. die Studienleistung mit „nicht bestanden“ bewertet wird
und bei mehrfachem oder schwerwiegendem Täuschungsversuch eine
Exmatrikulation erfolgen bzw. ein Verfahren zur Entziehung eines eventuell
verliehenen akademischen Titels eingeleitet werden kann.

\vspace*{1cm}

\hspace*{1cm}Tübingen, den 22. Mai 2026
\hfill 
Carlo Lanzi Luciani\hspace*{1cm}

\vspace*{1cm}







